% ============================================================
% Chapter 2 -- Background and related work
%
% The context chapter. It exists to make one distinction available to the rest
% of the thesis: the exact score, the empirical score and the fitted score are
% three different objects, and the difficulties of generative diffusion are
% properties of the gaps between them. Derivations of the textbook material
% summarised here are in the companion compendium.
% ============================================================

\chapter{Statistical Mechanics, Energy-Based Models, and Diffusion}
\label{ch:statmech}

This chapter reads generative diffusion through the statistical-mechanics
lineage it came from, fixes the score--posterior identity that every later
chapter uses, and locates the gap the thesis addresses; textbook material is
stated and cited, not derived.

\section{From the Boltzmann law to energy-based models}
\label{sec:sm-what}\label{sec:sm-ebm}

Statistical mechanics abandons the trajectory of a many-body system and keeps
a probability law over its microstates. Among all laws with a given mean
energy, the equilibrium one maximises the entropy
$S[p] = -\sum_x p(x)\log p(x)$ \citep{jaynes1957information}, and the
maximiser is the Boltzmann distribution
\begin{equation}
  p(x) \;=\; \frac{e^{-\beta E(x)}}{Z(\beta)},
  \qquad
  Z(\beta) \;=\; \sum_{x} e^{-\beta E(x)},
  \label{eq:sm-boltzmann}
\end{equation}
with the inverse temperature $\beta$ as the Lagrange multiplier of the energy
constraint \citep{landau1980statistical}. Two features of
\eqref{eq:sm-boltzmann} recur throughout. The macroscopic behaviour is encoded
in the single scalar $\log Z$, a sum over exponentially many configurations
that is $\#$P-hard for broad classes of interacting models
\citep{barahona1982computational, jerrum1993polynomial}; the tractable
exceptions, trees and Gaussians, are exactly the structures later chapters
exploit. And the competition between energy and entropy can produce
\emph{phase transitions}, changes of character of the equilibrium state as a
parameter is varied without any change in the microscopic law
\citep{ising1925beitrag, onsager1944crystal}.

Energy-based models read \eqref{eq:sm-boltzmann} in the opposite direction,
in a lineage that runs from the Hopfield network and its spin-glass capacity
analysis \citep{hopfield1982neural, amit1985storing, mezard1987spin} to the
Boltzmann machine \citep{ackley1985learning} and the 2024 Nobel Prize in
Physics \citep{nobel2024physics}: one posits a parametric energy $E_\theta$
and fits it so that $p_\theta(x) = e^{-E_\theta(x)}/Z(\theta)$ matches data
\citep{lecun2006tutorial}. The maximum-likelihood gradient is a competition
between two phases,
\begin{equation}
  \nabla_\theta\, \E_{x \sim p_{\mathrm{data}}}\!\left[-\log p_\theta(x)\right]
  \;=\;
  \E_{x^{+} \sim p_{\mathrm{data}}}\!\left[\nabla_\theta E_\theta(x^{+})\right]
  \;-\;
  \E_{x^{-} \sim p_{\theta}}\!\left[\nabla_\theta E_\theta(x^{-})\right],
  \label{eq:sm-ebm-grad}
\end{equation}
which follows from $\nabla_\theta \log Z(\theta) = -\E_{p_\theta}[\nabla_\theta
E_\theta]$: the positive phase lowers the energy of data, the negative phase
raises the energy of the model's own samples. The second expectation is over
$p_\theta$, which cannot be sampled directly; approximating it by Markov chain
Monte Carlo \citep{hinton2002training, robert2004monte} is what made the
family expensive and, for a long time, outperformed by models that give up the
explicit density for cheap sampling \citep{goodfellow2014generative,
kingma2014auto, papamakarios2021normalizing}. The training dynamics of these
models is itself a statistical-mechanics problem: \citet{bachtis2024cascade}
show that a restricted Boltzmann machine learns its features through a cascade
of second-order transitions, coarsest first, an ordering that
Section~\ref{sec:sm-transitions} meets again along the generation trajectory.

\section{Generative diffusion}
\label{sec:sm-diffusion}\label{sec:sm-relaxation}

A diffusion model destroys structure in a controlled way and learns to run the
destruction backwards \citep{sohl2015deep, song2019generative, ho2020denoising,
song2021scorebased}; its deterministic counterparts, neural ODEs and flow
matching, transport the same marginals along a deterministic path
\citep{chen2018neural, lipman2023flow}. The forward process is the SDE
\begin{equation}
  \dd X_t \;=\; f(X_t, t)\,\dd t \;+\; g(t)\,\dd W_t ,
  \qquad X_0 \sim p_0 ,
  \label{eq:sm-forward}
\end{equation}
whose marginals $p_t$ obey the Fokker--Planck equation
$\partial_t p_t = -\nabla\!\cdot(f p_t) + \tfrac12 g^2 \Delta p_t$
\citep{gardiner2009stochastic}. The choice used throughout this thesis is not
arbitrary: overdamped Langevin dynamics in a confining potential $E$ has the
Boltzmann law $\propto e^{-\beta E}$ as its unique stationary density, and
$\beta = 1$, $E(x) = \tfrac12\|x\|^2$ gives $f = -x$, $g = \sqrt2$, the
Ornstein--Uhlenbeck (OU) channel \citep{uhlenbeck1930theory}, whose transition
law is available in closed form,
\begin{equation}
  x \;=\; e^{-t} a \;+\; \sqrt{\Deltat}\,\varepsilon,
  \qquad \varepsilon \sim \Nd(0, I),
  \qquad \Deltat = 1 - e^{-2t}.
  \label{eq:sm-ou}
\end{equation}
The signal contracts by $e^{-t}$ while isotropic noise of variance $\Deltat$
fills in; the terminal law is the equilibrium of a quadratic Hamiltonian at
unit temperature, and the forward process is a relaxation to it from the data
law (Figure~\ref{fig:sm-corruption}). The original construction of
\citet{sohl2015deep} took exactly this non-equilibrium viewpoint: rather than
normalising a global energy, it compares forward and learned reverse
\emph{path} measures, and the resulting evidence lower bound is tractable
because every forward transition is known and locally normalised.

\begin{figure}[t]\centering
\includegraphics[width=\linewidth]{fig_forward_corruption.pdf}
\caption[The forward channel applied to a correlated sequence]{One clean
stationary AR(1) trajectory with $\alpha=0.85$ (grey) and the same trajectory
after the OU channel \eqref{eq:sm-ou} at four diffusion times. The signal
contracts by $e^{-t}$ while independent noise of variance $\Deltat$
accumulates; correlation between neighbouring sites outlives the original
amplitude. Generation runs this evolution from right to left.}
\label{fig:sm-corruption}
\end{figure}

What determines the reverse dynamics is a classical result. For a
state-independent $g(t)$ and smooth, positive marginals with integrable score,
\citet{anderson1982reverse} shows that
\begin{equation}
  \dd X_t
  \;=\;
  \big[ f(X_t,t) - g(t)^2\,\nabla_x\log p_t(X_t) \big]\dd t
  \;+\; g(t)\,\dd\bar W_t ,
  \label{eq:sm-reverse}
\end{equation}
integrated from $T$ back to $0$ from $X_T \sim p_T$, reproduces the forward
marginals in reverse order; the probability-flow ODE
$\dot X_t = f - \tfrac12 g^2 \nabla_x \log p_t$ has the same one-time
marginals with no Brownian term \citep{song2021scorebased}. Both are driven
by the same object, the \emph{score} of the noised law,
\begin{equation}
  s_t(x) \;:=\; \nabla_x \log p_t(x),
  \label{eq:intro-score}
\end{equation}
and so is every practical parameterisation: predicting the injected noise, as
in \citet{ho2020denoising}, is the same conditional expectation rescaled
\citep{lai2025principles}.

\section{The score--posterior identity, and three scores}
\label{sec:sm-three-scores}\label{sec:tweedie}

The word ``score'' is used for three different objects, and almost every
difficulty in generative diffusion is a property of the gaps between them
\citep{biroli2023generative}.

\paragraph{The exact score.} The score of the true noised law
$p_t = $ (contracted $p_0$) $* \Nd(0,\Deltat I)$. Writing $p_t(x) = \int
p_0(a)\, G(x; e^{-t}a, \Deltat)\,\dd a$ with $G$ the Gaussian kernel of
\eqref{eq:sm-ou}, differentiating under the integral and using
$\nabla_x G = -\Deltat^{-1}(x - e^{-t}a)\,G$ gives
\begin{equation}
  \boxed{\;
  \nabla_x \log p_t(x)
  \;=\; \frac{e^{-t}\,\E[\,a \mid x\,] - x}{\Deltat}
  \quad\Longleftrightarrow\quad
  \E[\,a \mid x\,]
  = e^{t}\big( x + \Deltat\, \nabla_x \log p_t(x) \big),
  \;}
  \label{eq:tweedie}
\end{equation}
where $\E[a \mid x]$ is the mean of the posterior
$p_t(a \mid x) \propto p_0(a)\,G(x; e^{-t}a, \Deltat)$, whose normaliser is
$p_t(x)$ itself. The identity needs only $\int \|a\|\,p_0(a)\,\dd a < \infty$.
It is Tweedie's formula of empirical Bayes \citep{robbins1956empirical,
miyasawa1961empirical, efron2011tweedie}, and it says more than an analogy:
computing a score \emph{is} solving a Bayesian denoising problem. For
vector-valued data $a \in \R^{\nsites}$ noised coordinatewise it holds
componentwise but with the \emph{full} posterior,
\begin{equation}
  \big( \nabla_x \log p_t(x) \big)_k
  \;=\; \frac{e^{-t}\, \E[\,a_k \mid x_0, \dots, x_{\nsites-1}\,] - x_k}{\Deltat},
  \label{eq:tweedie-joint}
\end{equation}
because the channel is independent across frames but the prior is not. This
is the exact sense in which the joint score of a trajectory differs from a
stack of per-frame scores, and the reason this thesis studies the former.

\paragraph{The empirical score.} Replacing $p_0$ by the empirical measure of a
training set $\{a^{(i)}\}_{i=1}^{N}$ makes $p_t$ a finite Gaussian mixture and
the score a softmax-weighted average of training points,
\begin{equation}
  \nabla_x \log p^{\mathrm{emp}}_t(x)
  \;=\; \frac{1}{\Deltat}\Big(
  e^{-t} \sum_{i=1}^{N} w_i(x)\, a^{(i)} - x \Big),
  \qquad
  w_i(x) \;\propto\; e^{-\lVert x - e^{-t} a^{(i)}\rVert^2 / 2\Deltat}.
  \label{eq:sm-empirical}
\end{equation}
It is exactly computable and exactly useless as a generative model: as
$t \to 0$ the weights concentrate on the nearest training example and the
reverse dynamics deposit the trajectory on top of it.

\paragraph{The fitted score.} A network $s_\phi$ trained by denoising score
matching \citep{hyvarinen2005estimation, vincent2011connection} minimises
\begin{equation}
  \E_{t}\, \E_{a \sim p_{\mathrm{data}}}\, \E_{\varepsilon}
  \Big\lVert
    s_\phi\big(e^{-t} a + \sqrt{\Deltat}\,\varepsilon,\; t\big)
    + \frac{\varepsilon}{\sqrt{\Deltat}}
  \Big\rVert^{2}.
  \label{eq:sm-dsm}
\end{equation}
At the population level its minimiser over all measurable functions is the
exact score; conditional on a finite training set it is the empirical score
\eqref{eq:sm-empirical}. A trained network reaches neither, and one proposed
mechanism of generalisation is precisely this gap: the network is too
constrained to represent \eqref{eq:sm-empirical}, and that failure is what
produces novel samples \citep{kadkhodaie2024generalization,
kamb2024analytic}. On that account generalisation is an
approximation-theoretic property of an architecture, a considerably less
comfortable foundation than the variational principle behind
\eqref{eq:sm-boltzmann}. Deciding what a trained network has learned therefore
requires the score it should have produced, and for essentially every data law
of interest that reference is unavailable. The solvable models of this thesis
supply it.

\section{Transitions along the generation trajectory}
\label{sec:sm-transitions}

Because generation is a dynamical process, its failures are dated by two
crossover times. \emph{Speciation}: run backwards from noise, the trajectory
commits to a broad class of outcomes long before it commits to a sample. In
the Gaussian-mixture models of \citet{biroli2024dynamical} and
\citet{achilli2026speciation}, for data whose covariance has top eigenvalue
$\lambda$, commitment happens near $t_{\mathrm S} \sim \tfrac12\log\lambda$,
when the contracted signal along that direction becomes comparable with the
accumulated noise; for hierarchical data there is a ladder of such times, one
per level, and features are resolved coarse to fine \citep{sclocchi2025phase,
sclocchi2025probing, favero2025compositional, favero2025bigger}, the ordering
\citet{bachtis2024cascade} observe in training. \emph{Collapse and
memorisation}: later in the reverse trajectory, dynamics driven by the
empirical score are captured by a single training example past a time
$t_{\mathrm C}$, and avoiding this needs a training set exponentially large in
the data dimension \citep{biroli2024dynamical, bonnaire2025memorize}; the
criterion is entropic, collapse occurring when the excess entropy of the
noised law over the terminal measure falls below $\log N$. Memorisation need
not be a collapse onto one example \citep{achilli2024losing}, the resulting
generalisation can be biased \citep{garnierbrun2026biased}, and solvable
models of the learning problem itself complete the picture
\citep{cui2025solvable}. The two transitions bracket the useful part of the
trajectory. The solvable-model analyses in this line treat a data point as an
unstructured vector; the sequence axis of this thesis is absent from them.

\section{Locality, architecture, and message passing}
\label{sec:sm-unet}

If the fitted score is doing the work, the question is what the architecture
contributes. The network in almost every image or video diffusion model is a
U-Net \citep{ronneberger2015unet}: a convolutional encoder that downsamples, a
decoder that upsamples, and long skip connections joining stages of equal
resolution. \citet{mei2024unets} gives that design a precise reading. For data
generated by a tree-structured hierarchical model the exact denoiser is
computed by belief propagation, an upward pass that accumulates evidence from
the leaves and a downward pass that redistributes it, reusing the upward
messages; a U-Net has exactly this shape, encoder, decoder and skips, so for
those models it is an approximate implementation of the exact inference
algorithm, with sample-complexity bounds to match. Nothing in that reading
certifies an arbitrary U-Net on arbitrary data as message passing, and the
locality that convolutional architectures build in \citep{lecun2015deep} is a
statement about receptive fields, not about the probabilistic structure of the
data.

Two consequences shape this thesis. Sum--product is exact on trees
(Theorem~\ref{thm:tree-exact}), so on a chain the posterior marginals, and by
\eqref{eq:tweedie-joint} the joint score, are computed with no approximation
at the inference stage: the structured reference against which learned
denoisers are compared. And the same formulation makes the receptive field a
quantity rather than a hyperparameter: the influence of an observation $d$
sites away decays at a rate the model fixes, which Chapter~\ref{ch:em-results}
computes in closed form and uses to predict the best window radius before any
network is trained. On sequences, \citet{garnierbrun2024transformers} use
exact belief propagation in the same role, as the reference for a
transformer trained on hierarchically generated data, and find that
correlations are learned shortest scale first.

\section{Positioning and gap}
\label{sec:sm-taxonomy}

Every generative family turns the same goal, drawing from a law $p_0$ nobody
can write down, into a map from a trivially sampled reference law to samples:
adversarial networks by a push-forward trained in a minimax game
\citep{goodfellow2014generative}, variational autoencoders by a stochastic
decoder trained on a likelihood bound \citep{kingma2014auto}, normalizing
flows by an invertible map with exact likelihood
\citep{papamakarios2021normalizing}, autoregressive models one coordinate at
a time \citep{vaswani2017attention}, and diffusion by repeated denoising,
sequential in the noise level rather than in the dimension;
\citet{ghio2024sampling} place flows, diffusion and autoregressive networks
under one spin-glass analysis and find the same first-order transition
governing sample quality in all three.

Three literatures meet in this thesis and leave a gap between them. The
statistical physics of diffusion has closed-form analyses of the backward
dynamics and its transitions, but for data without an internal time axis.
Exact inference on Gauss--Markov chains is classical, through the Kalman
filter and smoother and their reading as sum--product on a tree
\citep{kalman1960new, rauch1965maximum, pearl1988probabilistic,
kschischang2001factor, mezard2009information}, but it does not ask how the
exact posterior, and so the exact score, deforms as the noise level is swept
from zero to saturation. And message passing with continuous variables closes
its messages inside a parametric family, the step connecting belief
propagation to AMP and to the TAP equations \citep{thouless1977solution,
donoho2009message, zdeborova2016statistical}, on a central-limit argument that
needs many messages per node and fails by construction on a chain. What
follows computes the joint score of a noised Markov sequence exactly where
that is possible, measures the price of Gaussian closure where it is not,
and learns the dynamics from noised data when they are unknown, against a
reference that is deterministic quadrature rather than Monte Carlo, because a
discretisation error can be bounded separately and reported whereas a
sampling error would enter every comparison.
